\subsection{热统续}

\begin{结论}{Einstein 固体与自由度冻结}
$3N$ 个频率 $\omega_E$ 的量子振子，$x=\hbar\omega_E/kT$：
\[
\重点{C_V=3Nk\frac{x^2\ee^x}{(\ee^x-1)^2}},\qquad T\gg\Theta_E:\ 3Nk;\qquad T\ll\Theta_E:\ \propto x^2\ee^{-x}.
\]
能隙 $\Delta$ 使对应自由度在 $kT\ll\Delta$ 时指数冻结；能均分只对经典二次自由度成立（每个 $\frac12kT$）。低温实测为 $T^3$，说明需要 Debye 的连续声子谱。
\end{结论}

\begin{结论}{巨正则系综与粒子数涨落}
$\mathcal Z=\sum\ee^{-\beta(E-\mu N)}$，$\Omega=-kT\ln\mathcal Z=-pV$。单粒子态平均占据：
\[
\重点{\bar n=\frac{1}{\ee^{\beta(\varepsilon-\mu)}\pm1}}\ (+\text{ Fermi},\ -\text{ Bose}),\qquad \langle\Delta N^2\rangle=kT\Bigl(\frac{\partial\langle N\rangle}{\partial\mu}\Bigr)_{T,V}.
\]
单态涨落：Fermi $\bar n(1-\bar n)$，Bose $\bar n(1+\bar n)$（光子聚束），经典 $\bar n$（Poisson）。
\end{结论}

\begin{结论}{经典理想气体：化学势与 Sackur--Tetrode}
热波长 $\lambda_T=\dfrac{h}{\sqrt{2\pi mkT}}$：
\[
\重点{\mu=kT\ln(n\lambda_T^3)},\qquad S=Nk\Bigl[\ln\frac{V}{N\lambda_T^3}+\frac52\Bigr].
\]
经典条件即 $n\lambda_T^3\ll1$（此时 $\mu\ll0$）；$n\lambda_T^3\sim1$ 时量子简并（BEC、电子气）。$N!$ 保证熵广延（Gibbs 佯谬）。
\end{结论}

\begin{结论}{Saha 电离方程}
电离平衡 $\mathrm H\rightleftharpoons p+e$，电离能 $I$：
\[
\重点{\frac{n_en_p}{n_{\mathrm H}}=\Bigl(\frac{2\pi m_ekT}{h^2}\Bigr)^{3/2}\ee^{-I/kT}}\quad(\text{氢；一般情形乘 }2g_i/g_0).
\]
右边前因子很大，故 $kT\ll I$ 时也能大量电离：宇宙复合发生在 $T\approx3000\,$K（$kT\approx0.26\,$eV，远小于 $13.6\,$eV），因为光子--重子比 $\sim10^9$。
\end{结论}

\begin{结论}{Clausius--Clapeyron}
两相共存线：
\[
\重点{\frac{\dd p}{\dd T}=\frac{L}{T\,\Delta v}},\qquad \text{蒸气视为理想气体且 }L\text{ 恒定：}\ p\propto\ee^{-L/RT}.
\]
冰的 $\Delta v<0$，熔点随压强降低（约 $-0.0074\,$K/atm）。水的沸点每升高 $1\,$km 海拔约降 $3.3\,$K。
\end{结论}

\begin{结论}{热机与制冷的上限}
\[
\eta\le1-\frac{T_c}{T_h},\qquad \text{制冷 }\mathrm{COP}\le\frac{T_c}{T_h-T_c},\qquad \text{热泵 }\mathrm{COP}\le\frac{T_h}{T_h-T_c}.
\]
熵变只看初末态：理想气体自由膨胀到 $2V$，\重点{\Delta S=nR\ln2}（不可逆、$Q=0$）；两种气体等温等压混合 $\Delta S=-nR\sum x_i\ln x_i$。最大功 $=-\Delta F$（等温）。
\end{结论}

\begin{结论}{van der Waals 方程的临界点}
$\bigl(p+\frac{a}{V^2}\bigr)(V-b)=RT$（摩尔量）：
\[
\重点{T_c=\frac{8a}{27Rb},\quad p_c=\frac{a}{27b^2},\quad V_c=3b},\qquad Z_c=\frac{p_cV_c}{RT_c}=\frac38 .
\]
约化变量下方程普适（对应态定律）；$T<T_c$ 时用 Maxwell 等面积法则确定共存压强；临界指数为平均场值（$\beta=\frac12$）。
\end{结论}

\begin{结论}{绝热过程与大气递减率}
$pV^\gamma$、$TV^{\gamma-1}$ 恒定，做功 $W=\frac{p_1V_1-p_2V_2}{\gamma-1}$。绝热大气：
\[
\重点{\frac{\dd T}{\dd z}=-\frac{g}{c_p}\approx-9.8\,\mathrm{K/km}}\quad(\text{干空气}).
\]
湿空气凝结放潜热，约 $-6.5\,$K/km。实际递减率小于绝热值时大气对流稳定；等温大气标高 $H=\frac{kT}{mg}\approx8.4\,$km。
\end{结论}

\begin{结论}{密度涨落与压缩率}
\[
\重点{\frac{\langle\Delta N^2\rangle}{\langle N\rangle^2}=\frac{kT\kappa_T}{V}}.
\]
理想气体 $\kappa_T=1/p$，给出 Poisson 涨落 $1/\langle N\rangle$。临界点 $\kappa_T\to\infty$，涨落在光波长尺度上变大：临界乳光；光散射强度 $\propto\kappa_T$。
\end{结论}

\begin{结论}{Ising 模型：一维转移矩阵与二维 Onsager}
一维 $N$ 格周期边界：$Z=\lambda_+^N+\lambda_-^N$，
\[
\lambda_\pm=\ee^{\beta J}\cosh\beta h\pm\sqrt{\ee^{2\beta J}\sinh^2\beta h+\ee^{-2\beta J}},\qquad h=0:\ \xi^{-1}=\ln\coth\beta J .
\]
$\xi$ 只在 $T\to0$ 发散：一维无有限温相变。二维方格（Onsager）：\重点{\sinh\frac{2J}{kT_c}=1}，$kT_c=\frac{2J}{\ln(1+\sqrt2)}\approx2.269J$，$\beta=\frac18$，$\alpha=0$（对数发散）。
\end{结论}

\begin{结论}{顺磁体：$\tanh$ 与 Curie 定律}
$N$ 个自旋 $\frac12$、磁矩 $\mu$：
\[
\重点{M=N\mu\tanh\frac{\mu B}{kT}},\qquad \text{弱场 }M\approx\frac{N\mu^2B}{kT}\ (\text{Curie},\ \chi\propto1/T).
\]
一般角动量 $J$：$M\propto B_J(x)$（Brillouin 函数），Curie 常数 $\propto g^2J(J+1)\mu_B^2$。二能级布居反转对应负温度，熵对能量的导数为负。
\end{结论}

\begin{结论}{分子运动论：速率、自由程与粘滞}
Maxwell 分布：$v_p=\sqrt{\frac{2kT}{m}}$，\重点{\bar v=\sqrt{\frac{8kT}{\pi m}}}，$v_{\mathrm{rms}}=\sqrt{\frac{3kT}{m}}$；泻流通量 $\frac14n\bar v$。
平均自由程 $\lambda=\frac{1}{\sqrt2\,n\sigma}$；粘滞 $\eta\approx\frac13nm\bar v\lambda\propto\sqrt{mT}/\sigma$，\textbf{与压强无关}（Maxwell 的意外预言）；热导 $\kappa\approx\frac13nc_v\bar v\lambda$。
\end{结论}

\begin{结论}{Sommerfeld 展开}
$f$ 为 Fermi 函数，$kT\ll\mu$：
\[
\重点{\int_0^\infty H(\varepsilon)f(\varepsilon)\,\dd\varepsilon=\int_0^\mu H\,\dd\varepsilon+\frac{\pi^2}{6}(kT)^2H'(\mu)+O(T^4)}.
\]
三维自由电子：$\mu(T)=E_F\bigl[1-\frac{\pi^2}{12}(kT/E_F)^2\bigr]$；电子比热 $C=\frac{\pi^2}{3}g(E_F)k^2T$。热激发只发生在 $E_F$ 附近 $\sim kT$ 的薄层中。
\end{结论}

\begin{结论}{信息的热力学代价}
Landauer：擦除 1 比特至少耗散 \重点{kT\ln2}（$300\,$K 时 $2.9\times10^{-21}\,$J）。Szilard 机：知道分子在哪一半可提取 $kT\ln2$ 的功，记忆的擦除补回熵。
非平衡功涨落：Jarzynski $\langle\ee^{-\beta W}\rangle=\ee^{-\beta\Delta F}$，由 Jensen 得 $\langle W\rangle\ge\Delta F$；用于从单分子拉伸实验求自由能。
\end{结论}

\参考{Reif, \emph{Fundamentals of Statistical and Thermal Physics}；Kittel--Kroemer, \emph{Thermal Physics}；Pathria--Beale, \emph{Statistical Mechanics}；Kardar, \emph{Statistical Physics of Particles / Fields}（MIT 8.044 / 8.333 / 8.334）；Huang, \emph{Statistical Mechanics}。}
